\section{Étape 3 : calcul MPI complet avec cellules fantômes}

\subsection{Objectif}

L'objectif de cette étape est de distribuer réellement le calcul sur plusieurs processus MPI. Le \texttt{rank 0} reste dédié à l'affichage, mais les \texttt{ranks 1..P-1} se partagent maintenant la grille de calcul et ne conservent chacun qu'une partie des particules.

\subsection{Changements clés d'implémentation}

\begin{itemize}
    \item le \texttt{rank 0} lit le fichier initial une seule fois, crée un \texttt{global\_id} fixe pour chaque étoile et distribue les particules aux workers ;
    \item les workers forment une grille 2D logique sur le plan \(Oxy\), avec une décomposition en blocs rectangulaires de cellules ;
    \item chaque worker stocke ses particules propriétaires, plus des copies fantômes nécessaires aux interactions directes proches ;
    \item après chaque pas, les positions sont réunies sur le \texttt{rank 0} par \texttt{global\_id} pour reconstruire l'ordre global avant affichage.
\end{itemize}

\subsection{Topologie retenue}

Conformément au cours sur la décomposition de domaine, on ne découpe pas l'espace en bandes 1D. Les workers sont organisés en une grille 2D \(p_x \times p_y\) sur \(Ox\) et \(Oy\), tandis que l'axe \(Oz\) reste entier sur chaque worker.

Cette étape n'accepte donc que les exécutions où le nombre de workers se factorise en une vraie grille 2D, avec \(p_x > 1\) et \(p_y > 1\). Le cas à deux processus reste celui de l'étape 2.

\subsection{Rôle des cellules fantômes et des réductions globales}

Le noyau gravitationnel utilise un calcul direct dans les cellules voisines situées à distance de Chebyshev inférieure ou égale à 2. C'est la raison pour laquelle le halo MPI est fixé à deux cellules dans \(x\) et \(y\), diagonales incluses.

Pour les cellules lointaines, chaque worker n'utilise pas les particules une à une, mais la masse totale et le centre de masse globaux de chaque cellule. Ces quantités sont reconstruites par \texttt{Allreduce} en sommant uniquement les particules propriétaires. Les particules fantômes ne servent donc qu'aux interactions proches et ne participent jamais au comptage global.
